\subsection{Greedy por Alternancia con Límite de Racha (dónde dormir cada noche)}
\label{alg:8-4}

\graybold{Preguntas guía:}

\begin{itemize}
\item ¿Hay que construir explícitamente una secuencia que respete un límite de repeticiones CONSECUTIVAS de cada elemento (sin límite en el total de usos)?
\item ¿Un elemento puede reutilizarse libremente después de "descansar" al menos un turno?
\item ¿Basta con dos "recursos" (el mejor y el segundo mejor) para romper cualquier racha, sin necesidad de planear la secuencia completa de antemano?
\end{itemize}

\graybold{Utilidad:} Construir una secuencia válida de longitud D usando N opciones con límites de racha consecutiva, sin necesidad de una búsqueda ni un algoritmo sofisticado: alternar entre las DOS opciones con mayor límite es siempre suficiente (si el problema es factible), porque cualquier racha se puede "romper" con un solo turno de la segunda opción antes de retomar la primera.

\graybold{Complejidad:} \bigO{D}.

\graybold{Fundamento teórico:} El problema es imposible si ningún lugar admite ni una sola noche (todos los límites son 0), o si solo UN lugar tiene límite positivo y ese límite es menor que D (no hay forma de "descansar" de él). En cualquier otro caso (al menos dos lugares con límite $>0$), siempre es posible: se usa el lugar de mayor límite hasta agotar su racha, se intercala un solo día en el segundo mejor lugar (lo cual reinicia la racha del primero), y se repite este patrón hasta completar los D días. Este patrón simple es suficiente porque no hay ningún límite en el número TOTAL de veces que se puede volver a un lugar, solo en la racha consecutiva.

\graybold{Ejercicio aplicado}

Kai necesita dormir D días seguidos, alternando entre N lugares, donde el lugar $i$ tolera a lo más $d_i$ noches CONSECUTIVAS antes de expulsarlo (pero puede volver después). Construir un plan válido día a día, o determinar que es imposible.

\graybold{I/O:} N,D ≤ \ten{5} → lectura total con tokenización (patrón 2). Verificado contra 200 casos aleatorios comparando factibilidad y validando la secuencia construida.

\graybold{Código solución}

\begin{lstlisting}[language=Python]
import sys

def solve():
    data = sys.stdin.buffer.read().split()
    n = int(data[0]); d = int(data[1])
    arr = list(map(int, data[2:2 + n]))

    order = sorted(range(n), key=lambda i: -arr[i])
    best_idx = order[0]
    best_val = arr[best_idx]

    if best_val == 0:                      # ningun lugar es utilizable
        print(-1)
        return

    second_val = arr[order[1]] if n > 1 else 0

    if second_val == 0:
        # solo hay UN lugar utilizable: hay que quedarse ahi todos los D dias seguidos
        if best_val >= d:
            print(' '.join([str(best_idx + 1)] * d))
        else:
            print(-1)
        return

    second_idx = order[1]
    out = []
    remaining = d
    while remaining > 0:
        if remaining == 1:
            out.append(second_idx + 1)          # un dia de "descanso" en el segundo mejor lugar
            break

        out.append(second_idx + 1)          # un dia de "descanso" en el segundo mejor lugar
        remaining -= 1

        take = min(best_val, remaining)
        out.extend([best_idx + 1] * take)
        remaining -= take
    sys.stdout.write(' '.join(map(str, out)) + '\n')

solve()
\end{lstlisting}
